\documentclass[fleqn,10pt]{wlscirep}
\usepackage[utf8]{inputenc}
\usepackage[T1]{fontenc}
\usepackage{amsmath}
\usepackage{amssymb}
\usepackage{booktabs}
\usepackage{threeparttable}
\usepackage{multirow}
\usepackage{rotating}

\providecommand{\botrule}{\bottomrule}

% Scientific Reports asks for inline citations in square brackets. wlscirep.cls
% loads cite.sty with the [superscript] option; this restores bracketed form.
\makeatletter
\renewcommand\@citess[1]{[#1]}
\makeatother

\title{PLEF: An Evaluation of a Lexicon-Based Psycholinguistic Framework for Relationship-Related Text}

\author[1,*]{Rahul Kumar Mishra}

\affil[1]{Independent Researcher, Ranchi, Jharkhand, India 834012. ORCID: \href{https://orcid.org/0009-0004-9671-4055}{0009-0004-9671-4055}}

\affil[*]{mishra.rahul98@gmail.com}

\keywords{Narrative sentiment, Lexicon methods, Null models, Construct validity, Annotation reliability, Reproducibility}

\begin{abstract}
Metrics that summarise the emotional trajectory of a narrative -- where its direction changes, how volatile it is, whether it rises or falls overall -- are used in narrative analysis and are sometimes reported without an explicit structural null. This paper evaluates the Psycholinguistic Lexical Extraction Framework (PLEF), a standard-library-only core implementation of twelve interpretable narrative metrics, on 309{,}833 corpus records spanning seven public corpora and a 1{,}308-post relationship corpus, which reduce to 265{,}452 distinct documents after within- and cross-corpus text deduplication; analyses retain corpus identity and use the per-corpus records. Two results are primary. First, the association between the framework's two trajectory metrics, which would ordinarily be read as convergent validity, is not unusually large relative to an algorithmically specified within-document order-randomisation null: on 1{,}288 relationship-related posts the observed correlation is $r=+0.598$ (95\% CI $[+0.547,+0.646]$), while 10{,}000 within-document permutations, which preserve every document's sentence-level values and destroy only their order, give a null mean of $+0.607$ (SD $0.026$) and an empirical $p=0.643$. The two metrics are pre-minus-post mean differences over one signal, and in this dataset the observed association does not exceed the association generated by the specified within-document order-randomisation null. Second, the watershed metric did not match the human criterion in this preliminary evaluation, on a sample too small to distinguish that from chance: among 35 unique posts for which both a human response and a defined watershed index were available, the annotator judged that 16 contained no turning point at all, while the metric returned a defined index on all 35 eligible posts; on the remaining 19 it matched the annotator's sentence on none, against an expected 0.89 matches under the support-aware chance model. In a separate three-annotator round there was no unanimous three-annotator agreement on any of 26 jointly labelled posts, and within-annotator repeat agreement was 18--29\%, so the turning-point judgement was not reliably answerable under the annotation protocol used here. Two distinct conclusions follow and should not be conflated: the metric did not match this annotator's responses, and the available human criterion is too unreliable to support a strong criterion-validity conclusion. Against reference implementations of five lexical sentiment systems the framework does not lead on any of the evaluated corpora, and its 348-entry lexicon matches 1.37--2.95\% of tokens, producing a non-zero score on 7.2--23.8\% of documents of 1--9 words and 62.5--86.3\% of documents of 80--159 words. The contribution is therefore methodological: a fully specified permutation null against which trajectory-metric correlations can be tested, a measured coverage limitation for this lexicon-based implementation, and an abstention-aware evaluation that separates absence of lexical evidence from a neutral prediction.
\end{abstract}

\begin{document}

\hypersetup{
  pdftitle={PLEF: An Evaluation of a Lexicon-Based Psycholinguistic Framework for Relationship-Related Text},
  pdfauthor={Rahul Kumar Mishra},
  pdflang={en-US}
}

\flushbottom
\maketitle
\thispagestyle{empty}

\section{Introduction}\label{sec:intro}

Text analysis of personal narratives increasingly relies on metrics that summarise not the overall sentiment of a document but the \emph{shape} of its emotional trajectory: where the direction changes, how much it fluctuates, whether it ends better or worse than it began. Reagan et al.\cite{reagan2016} showed that a large corpus of fiction is well described by a small number of such arcs, and the narrative-identity literature treats self-told narratives as psychologically meaningful\cite{mcadams2013}. The appeal in applied settings is obvious, and metrics of this kind are now applied to personal and social-media text.

Trajectory metrics have a property that distinguishes them from ordinary sentiment scores and that is easy to overlook. They are functions of the \emph{same} underlying per-sentence signal, and two such functions can be strongly correlated for purely algebraic reasons, with no contribution from the narrative structure they are intended to capture. A correlation between two such metrics is therefore not, by itself, sufficient evidence that either measures narrative shape, unless it is shown to exceed what the shared signal produces on its own. Such correlations can nonetheless be interpreted as convergent validity, and the null against which they should be judged is not always reported.

This paper reports what happens when that comparison is made. The vehicle is the Psycholinguistic Lexical Extraction Framework (PLEF), a small, deterministic implementation of twelve interpretable metrics designed for longer-form narrative text in the relationship domain. Two of its metrics are trajectory metrics: a watershed index (LEWI), defined as the point of maximum absolute second finite difference of the smoothed sentence-level sentiment signal, and an arc asymmetry (NAVA), the difference between the mean of the opening sentences and the mean of the closing ones. Both are pre-minus-post mean differences over one signal, differing only in where the split falls.

This paper reports a null for their association, reports the observed value against it, and reports a second, independently measured criterion test: whether the watershed index agrees with where a human reader says a narrative turns.

Three further questions follow from the corrected evaluation. What does a 348-entry lexicon actually see, and at what document length does it see enough to produce a signal? When a lexicon matches nothing and the framework returns zero, is that a prediction of neutrality or an abstention, and does the distinction change the measured accuracy? And is the human judgement that the watershed metric targets reliable enough to serve as a criterion at all?

The contributions are: (i) a within-document permutation null for trajectory-metric correlations, and the observation that a correlation of this kind does not exceed it; (ii) a direct test of a watershed metric against an independent human annotator, including the finding that the metric cannot represent the absence of a turning point; (iii) a reliability measurement showing that the turning-point judgement is not consistently reproducible between or within annotators; and (iv) a coverage and document-length characterisation of the empirical operating range of the instrument in these corpora.

\section{Methods}\label{sec:methods}

\subsection{The framework}\label{subsec:framework}

PLEF computes twelve metrics from a document. Sentence-level sentiment is a lexicon sum with negation and intensity modifiers, normalised as $C(s) = x(s)/\sqrt{x(s)^2+15}$, where $x(s)$ is the raw lexicon-weighted sentence score, following Hutto and Gilbert\cite{hutto2014}. The symbol $\sigma$ is reserved throughout for standard deviations, over a hand-built 348-entry lexicon oriented toward relationship discourse. The 348 entries were written on 18 April 2026, twenty-two days before the relationship corpus file used here was assembled on 10 May 2026. The lexicon therefore predates this corpus and was not derived from it, and no post from it, and no label of any kind, was consulted while the entries were chosen. These records support the statement that the lexicon was constructed without consulting the evaluation corpus or its labels; they do not and could not establish that the author was unfamiliar with relationship discourse in general, which is the domain the lexicon was written for. The lexicon was frozen at the start of the analysis reported here and was not modified during it; the released file \texttt{plef\_lexicons.json} has SHA-256 hash \texttt{18294c67\dots25165c} and is recorded in the release manifest. From the resulting per-sentence vector $S=(s_1,\dots,s_N)$ the framework derives the metrics used here.

The two trajectory metrics are defined as follows. Let $w=\max(1,\lfloor N/4 \rfloor)$ and let $\tilde S$ be the centred moving
average of $S$ with half-width $w$: $\tilde s_i$ is the mean of
$s_{\max(1,\,i-w)}, \dots, s_{\min(N,\,i+w)}$. The window is truncated at the
boundaries and is never padded or reflected. The watershed index is
\begin{equation}
i^{*} = \operatorname*{arg\,max}_{2 \le i \le N-1}
\left| \tilde s_{i+1} - 2\tilde s_i + \tilde s_{i-1} \right|,
\qquad
\Delta_{\mathrm{LEWI}} = \frac{1}{i^{*}-1}\sum_{j=1}^{i^{*}-1} s_j
\; - \; \frac{1}{N-i^{*}+1}\sum_{j=i^{*}}^{N} s_j ,
\label{eq:lewi}
\end{equation}
with the two segments disjoint and exhaustive. The search is restricted to $2 \le i \le N-1$ so that the second finite difference, which requires positions $i-1$, $i$ and $i+1$, is defined. With this restriction the earlier segment contains $i^{*}-1 \ge 1$ sentences and the later one $N-i^{*}+1 \ge 2$.
When several interior indices attain the same maximum absolute second
difference, the implementation selects the smallest such index: it scans
interior positions in increasing order and replaces the incumbent only on a
strictly larger value. When $\tilde s$ is constant, so that no interior index
attains a strictly larger value than the initial candidate, it takes
$i^{*} = \lfloor N/2 \rfloor + 1$. Both rules are stated because the watershed
is sensitive to its own definition.
The arc asymmetry compares the opening and closing portions of the sequence.
Its window is
\begin{equation}
k = \begin{cases}
\max\!\left(1,\ \lfloor N/4 \rfloor\right), & N < 9,\\[2pt]
\max\!\left(1,\ \lfloor N/3 \rfloor\right), & N \ge 9,
\end{cases}
\qquad
\mathrm{NAVA} = \frac{1}{k}\sum_{j=1}^{k} s_j \; - \; \frac{1}{k}\sum_{j=N-k+1}^{N} s_j .
\label{eq:nava}
\end{equation}
Here $k$ is the arc window and is distinct from the smoothing half-width $w$ of Equation~\eqref{eq:lewi}: $w$ is $\max(1,\lfloor N/4 \rfloor)$ for every $N$, whereas $k$ switches to $\lfloor N/3 \rfloor$ at $N \ge 9$, so the two coincide only for $4 \le N \le 8$. The piecewise form is a property of the implementation as written rather than a derived optimum, and the threshold at nine sentences is implementation-specific rather than theoretically motivated; the primary analysis is therefore repeated under two fixed-window alternatives. Because the arc metric is one half of the primary null analysis, that analysis was repeated under three window rules: $k=\lfloor N/4 \rfloor$ throughout, $k=\lfloor N/3 \rfloor$ throughout, and the piecewise rule. The inferential conclusion is unchanged across all three specifications: none yields an observed correlation unusually large relative to its corresponding null. The values are given in Section~\ref{subsec:null} and the full comparison in the supplementary material. The two windows are disjoint for every $N \ge 4$: for $N<9$, $2k \le N/2$, and for $N \ge 9$, $2k \le 2N/3$, so no sentence contributes to both means. The tightest case is $N=9$, where the windows cover six of the nine sentences. Positive NAVA means the narrative
ends worse than it begins. Both are undefined for $N<4$ and are returned as missing rather than as zero, so that an undefined value cannot enter a correlation as a real one.

Equations~\eqref{eq:lewi} and~\eqref{eq:nava} are both differences between the mean of an early portion of $S$ and the mean of a late portion, differing only in where the split falls. This is the algebraic relationship that Section~\ref{subsec:null} tests.

The volatility metric is the mean absolute successive difference, $\mathrm{TEG} = (N-1)^{-1}\sum_{i=2}^{N}|s_i - s_{i-1}|$. This statistic is not novel: related successive-difference and inertia measures are used in affect-dynamics research to characterise short-term emotion change\cite{kuppens2010,houben2015}; the statistic is applied here to sentence-level narrative sentiment.

Two pronoun-derived indices require care. A pronoun triangulation index is defined as $(I-\mathrm{You})/(I+\mathrm{You}+\mathrm{We}+1)$. A power asymmetry index defined so as to include a term $|I-\mathrm{You}|/(I+\mathrm{You})$, with that term set to zero when a document contains no first- or second-person pronoun so the quotient is defined everywhere, would share its inputs with the first, and any association between the two would then be partly an algebraic identity rather than an empirical finding. The power index used here therefore excludes that term and is the mean of a control-verb density and an apology asymmetry, neither derived from pronoun counts. Section~\ref{subsec:remaining} reports the association under both definitions so that the size of the coupling is visible.

The twelve metric families are, in order: (1) mean sentence-level compound sentiment; (2) the watershed index LEWI and its drop; (3) the arc asymmetry NAVA with its five-class label; (4) the volatility index TEG; (5) the pronoun triangulation index PTI; (6) the power asymmetry index PAI; (7) the relational coherence index RCI; (8) the cognitive distortion index; (9) the vulnerability-disclosure score VADS; (10) the inconsistency index TIES; (11) a latent-semantic coherence measure; and (12) the composite relationship-health metric GASE. The watershed drop and the arc's five-class label are derived outputs of families (2) and (3) rather than separate families, so the number of reported quantities exceeds twelve. A further quantity, a semantic-fusion score combining GASE, semantic coherence, TEG and PTI, is derived from these rather than computed independently and is reported only as an internal diagnostic. Complete definitions of every metric, including normalisation, edge cases and missing-value handling, are given in the supplementary material; the released core module is the authoritative specification.

The semantic layer is latent semantic analysis\cite{deerwester1990,landauer1998}: term-frequency--inverse-document-frequency weighting followed by a truncated singular value decomposition of three components, computed by deterministic power iteration from a fixed initial vector; the parameters are given in the supplementary material. It is a linear matrix factorisation and is not described here as neural. The implementation uses one deterministic code path. Reproducibility is established empirically rather than by argument: the output is tested for byte-for-byte identity across separate processes.

The framework core is the module implementing tokenisation, sentence splitting, the sentiment scorer, the twelve metrics and the latent-semantic layer, together with the lexicon file it reads at import. It imports only the Python standard library. The evaluation harness, which loads the reference baselines and produces the figures, does not, and the two are separate. The claim is that the core has no third-party code dependency; it is not a claim that the framework is data-free, since the core requires its lexicon file. This is verified by a test that imports it in a fresh process with third-party packages blocked. Output is deterministic: a test scores an identical input set in two separate processes and compares the serialised results byte for byte.

\subsection{Corpora}\label{subsec:corpora}

We use \emph{document} as the generic unit of analysis throughout; the corpus-specific source unit is stated for each corpus, since these differ. Seven public corpora and one collected corpus are used (Table~\ref{tab:corpora}): GoEmotions\cite{demszky2020}, EmpatheticDialogues\cite{rashkin2019}, ISEAR\cite{scherer1994}, SemEval-2017 Task 4A\cite{rosenthal2017}, MELD\cite{poria2019}, TweetEval\cite{barbieri2020} and DailyDialog\cite{li2017}, together with one collected corpus of relationship-related Reddit posts. Where labels are used they are the datasets' own annotations; the relationship corpus contributes no labels to the corpus-level benchmark comparison. A separate 36-post exploratory point-estimate comparison on its human-labelled subset is reported in Section~\ref{subsec:limitations}. ISEAR has no neutral category, and EmpatheticDialogues has no label that maps to neutral; both are therefore evaluated and reported as two-class problems and are never pooled with the three-class corpora. Four EmpatheticDialogues rows whose emotion field matched neither class were found on inspection to be CSV parse fragments rather than labels, and were counted and dropped rather than assigned to a class. The label mapping for every corpus, with the resulting class distribution, is tabulated in the supplementary material.

The relationship corpus consists of 1{,}308 public posts collected from seven public subreddits, selected purposively for relevance to relationship and emotional-support discourse rather than as a representative sample of any population: r/relationships, r/relationship\_advice, r/BreakUps, r/survivinginfidelity, r/confession, r/emotionalsupport and r/heartbreak. Posts were retrieved through the public Reddit API; the corpus file used in this study was created on 10 May 2026, which bounds the collection date from above. Post-level timestamps were not retained in the corpus, so neither the exact query date nor the date range of the posts themselves can be recovered, and both are reported as unavailable rather than reconstructed. only post bodies were retained, and comments, deleted and removed posts, and posts under 40 words were excluded. Exact-duplicate bodies were removed at collection. No explicit language filter was applied. The collection script and the resulting identifier list are released; post text is not redistributed. Beyond these preprocessing exclusions, no content-based eligibility criterion was applied; in particular, no criterion was applied to establish that the posts are first-person accounts. Because no explicit language filter was used, the language composition of the collected corpus was not independently verified, which could also affect lexical coverage. Their length distribution, in the framework's own tokenisation, is mean 519 words (SD 489, range 48--6{,}591, median 394.5) and mean 32 sentences (SD 33, median 24). We did not identify labels for it of sufficient reliability for the classification analysis. Keyword-derived labels can be constructed, but on the 36 posts that also carry a human sentiment label they agree with the annotator on 25.0\% of items, with Cohen's $\kappa$\cite{cohen1960} $=-0.028$; the contingency counts and the keyword rule, including the classes it can emit and its handling of documents matching no keyword, are given in the supplementary material. They are therefore not used. No corpus-level classification result is computed on this corpus; an exploratory development-set comparison on the subset carrying human labels is reported separately in Section~\ref{subsec:limitations}. It is used for the trajectory analyses, which require no labels, and a reserved 200-post annotation-development subset, drawn from the same 1{,}508-post collection before the 1{,}308-post evaluation set was formed and therefore disjoint from it, is used for the human annotation described in Section~\ref{subsec:annotation}. No model parameter was tuned on this subset, no human label enters the trajectory-null analysis, and no post in the annotation-development subset appears among the 1{,}288 trajectory-eligible posts.

SemEval-2017\cite{rosenthal2017} and TweetEval\cite{barbieri2020} share 22{,}603 items, 42.8\% of the former, because TweetEval's sentiment task is derived from SemEval. Any mean over both double-weights those items; where a pooled mean is reported, the effect of excluding TweetEval is reported alongside it.

The analytic population must be stated explicitly, because two counts appear in this paper and they serve different purposes. \emph{All benchmark and metric analyses operate on each corpus's own records, retaining corpus identity and therefore retaining cross-corpus duplicates}; the per-corpus item counts of Table~\ref{tab:corpora} are the denominators for every result reported here. Global deduplication is used only to report how many distinct documents the study touches in total, and never as an analysis population. The 309{,}833 records support the coverage and classification analyses; the trajectory conclusion rests on the 1{,}288 trajectory-eligible relationship posts alone. This is why the SemEval--TweetEval overlap can double-weight a mean over both corpora, and why that effect is reported separately.

Item counts reduce in two distinct steps, which Table~\ref{tab:reconcile} separates. The eight corpora contribute 309{,}833 scored documents. Removing documents whose normalised text is duplicated \emph{within} their own corpus leaves 288{,}110. Removing documents duplicated \emph{across} corpora leaves 265{,}452, a further 22{,}658 documents. The pairwise overlaps that produce this are 22{,}603 between SemEval-2017 and TweetEval, 45 between GoEmotions and MELD, and 11 distributed across seven further pairs, summing to 22{,}659. That sum exceeds the number of documents removed by one because a document present in $k$ corpora appears in $k(k-1)/2$ pairs but is removed $k-1$ times; an excess of one indicates a single document present in three corpora. Reporting a single reduction from 309{,}833 to 265{,}452 would conflate the two steps and make neither reconstructible.

\begin{table}[ht]
\centering
\small
\caption{Corpora. ``Usable for trajectory'' counts items with at least four sentences, the minimum at which the watershed and arc metrics are defined. Word and sentence counts use the framework's own tokenisation, the same definition used for the length bands in Section~\ref{subsec:coverage}, so that a single length definition is used throughout. Medians of even-sized samples are the mean of the two central values. The final row sums the item and trajectory columns; medians are not pooled, since a median of medians has no interpretation. Item deduplication is reported separately in Table~\ref{tab:reconcile}.}
\label{tab:corpora}
\begin{tabular}{lrrrrl}
\toprule
Corpus & Items & Med.\ words & Med.\ sent.\ & $\ge 4$ sentences & Labels \\
\midrule
relationship   & 1{,}308   & 394.5 & 24 & 1{,}288 (98.5\%) & none \\
goemotions     & 54{,}263  & 12  & 1  & 914 (1.7\%)      & 3-cl. \\
empathetic     & 19{,}157  & 16  & 2  & 290 (1.5\%)      & 2-cl. \\
isear          & 7{,}666   & 18  & 1  & 296 (3.9\%)      & 2-cl. \\
semeval2017    & 52{,}806  & 21  & 2  & 3{,}405 (6.4\%)  & 3-class \\
meld           & 13{,}708  & 7   & 1  & 354 (2.6\%)      & 3-class \\
tweeteval      & 59{,}899  & 19  & 2  & 2{,}584 (4.3\%)  & 3-class \\
dailydialog    & 101{,}026 & 9   & 1  & 3{,}892 (3.9\%)  & 3-class \\
\midrule
\midrule
Sum            & 309{,}833 & --- & --- & 13{,}023          & --- \\
\botrule
\end{tabular}
\end{table}

The last column of Table~\ref{tab:corpora} determines the scope of the trajectory analysis. On seven of eight corpora fewer than 7\% of items have four or more sentences, and that surviving fraction is a length-selected subset of an overwhelmingly short-form distribution rather than a simple random sample of it. Reporting a correlation over such a tail and presenting it as a property of the corpus is not defensible. The trajectory metrics are therefore reported on the relationship corpus only, where 98.5\% of items qualify; for every other corpus they are reported as not measurable, with the count stated.

\begin{table}[ht]
\centering
\small
\caption{Item reconciliation. ``Raw units'' are the records read from the source files; discards are itemised with their cause. The two deduplication steps are separate and are reported separately.}
\label{tab:reconcile}
\begin{threeparttable}
\begin{tabular}{lrrrl}
\toprule
Corpus & Raw units & Retained & Discarded & Cause of discard \\
\midrule
dailydialog  & 102{,}979 & 101{,}026 & 1{,}953 & unparseable label field \\
empathetic   & 19{,}161  & 19{,}157  & 4   & malformed emotion field \\
goemotions   & 54{,}263  & 54{,}263  & 0   & --- \\
isear        & 7{,}666   & 7{,}666   & 0   & --- \\
meld         & 13{,}708  & 13{,}708  & 0   & --- \\
relationship & 1{,}308   & 1{,}308   & 0   & --- \\
semeval2017  & 53{,}608  & 52{,}806  & 802 & duplicate ids (764), malformed (38) \\
tweeteval    & 59{,}899  & 59{,}899  & 0   & --- \\
\midrule
Scored documents & & 309{,}833 & & \\
After within-corpus deduplication & & 288{,}110 & 21{,}723 & repeated text within a corpus \\
After cross-corpus deduplication & & 265{,}452 & 22{,}658 & across corpora\tnote{a} \\
\botrule
\end{tabular}
\begin{tablenotes}\footnotesize
\item[a] Pairwise overlaps: 22{,}603 SemEval--TweetEval, 45 GoEmotions--MELD, 11 across seven further pairs. The pairwise total of 22{,}659 exceeds the 22{,}658 documents removed by one, because one document is present in three corpora and so appears in three pairs while being removed twice.
\end{tablenotes}
\end{threeparttable}
\end{table}

Two corpora require a note on their relation to the canonical releases. We use the agreement-filtered GoEmotions release of 54{,}263 examples (43{,}410 train, 5{,}426 validation, 5{,}427 test), not the full 58{,}009-example release. For EmpatheticDialogues the document is the \emph{situation description}: the single free-text prompt describing the speaker's circumstances, which carries the self-reported emotion label. Neither the individual utterances nor the concatenated dialogue is scored. After parsing and deduplication this yields 19{,}157 unique situation descriptions, rather than the approximately 25{,}000 conversations of the canonical description. For DailyDialog the document is the individual \emph{utterance}, which carries the emotion label; utterances are short, and only 3.9\% contain four or more sentences, which is why the trajectory metrics are not computed on that corpus. These counts result from the specific source release and document unit used in this study; neither represents the full canonical release, and Table~\ref{tab:reconcile} gives the complete audit trail.

\subsection{Baselines}\label{subsec:baselines}

Five reference lexical sentiment systems are used for polarity classification: VADER\cite{hutto2014}, the NRC Emotion Lexicon\cite{mohammad2013}, Empath\cite{fast2016}, AFINN\cite{nielsen2011} and the Hu and Liu opinion lexicon\cite{hu2004}. labMT\cite{dodds2010} is considered separately as a happiness lexicon and is not used for polarity classification. LIWC\cite{tausczik2010,boyd2022} is not used: it is proprietary and no licence was held. Empath\cite{fast2016} is included instead as an openly available lexical-category comparator. The word LIWC is not used to name a baseline anywhere in this paper.

labMT is excluded from the polarity comparison because it measures happiness rather than polarity and its zero point is not neutral: on relationship text it labels 97.8\% of documents positive, with Cohen's $\kappa$ between $0.013$ and $0.039$ against every other system. It appears in no results table in this paper and no conclusion rests on it.

Each baseline is converted to a single signed document score before thresholding, and the conversions differ because the lexicons differ in kind. VADER contributes its own compound score. AFINN defines token-level valences but prescribes no document-level normalisation for a classification task of this kind, so a study-specific adapter normalisation was required; the sum is divided by five times the square root of the document's total token count and clipped to $[-1,1]$; with no scored token the numerator is zero and the score is zero. The Hu and Liu lexicon, which is a bare positive/negative word list, gives $(p-n)/(p+n)$ over matched tokens, and zero when no token matches. The NRC Emotion Lexicon is an emotion-association resource rather than a polarity scorer; only its positive and negative association flags are used, in the same $(p-n)/(p+n)$ form with the same zero-match rule, and its eight emotion categories are ignored. Empath is a lexical-category instrument: its normalised category scores are summed over 19 positive and 26 negative categories, chosen by name, and combined as $(p-n)/(p+n)$, again zero when both sums vanish. The narrow positive-emotion and negative-emotion pair was tried first and returned exactly zero on two of five hand-written test sentences, after which the broader set was adopted; the final category set was therefore selected during adapter development, using five hand-written sentences that belong to no evaluation corpus; the adapter was then frozen before any evaluation corpus was scored, and both variants are scored in the released code. This is development-time tuning rather than a choice specified in advance, and it is classified as such in Table~\ref{tab:dev}. Every conversion operates on the same tokenised representation described in Section~\ref{subsec:corpora}, and the released baseline adapter implements these definitions. Two of the conversions are adapter-defined rather than native outputs of the resource: the NRC Emotion Lexicon is an emotion-association resource and Empath is a lexical-category instrument, so in both cases the signed score reported here is a conversion chosen to produce a quantity compatible with the common three-class evaluation, and should not be read as the resource's own polarity output. A single classification threshold of $\pm 0.05$ is applied to all systems: a score strictly greater than $+0.05$ is positive, strictly less than $-0.05$ is negative, and any score in the closed interval $[-0.05, +0.05]$, including a score of exactly zero, is neutral. A system that returns zero because nothing matched is therefore labelled neutral, which is the mapping Section~\ref{subsec:abstention} examines. This is VADER's convention and it is not natural for the others, so macro-F1 is treated as threshold-conditional and the mean binary AUC over positive-vs-rest and negative-vs-rest, which is threshold-free, is the primary metric. The ranking reported below rests on a common threshold rather than per-system threshold optimisation; the supplementary threshold sweep assesses sensitivity to that choice, and shows that the framework would not lead on any three-class corpus even if every system were evaluated at its own best threshold. A threshold sweep from 0.00 to 0.30 for every system on every corpus is tabulated in the supplementary material, together with a comparison at each system's own best threshold. For three-class corpora the primary threshold-free metric is the mean of the positive-versus-rest and negative-versus-rest one-versus-rest AUCs, referred to hereafter as the mean binary AUC over positive-vs-rest and negative-vs-rest: because neutral occupies the centre of the signed scale, a neutral-versus-rest ranking would require an additional transformation, which is not specified here. This limitation is stated explicitly because a three-class macro-AUC computed from a single scalar score is otherwise ambiguous.

\subsection{Nulls and intervals}\label{subsec:nullmethod}

All reported correlations are Pearson product-moment correlation coefficients. Confidence intervals for benchmark scores and for correlations are bootstrap intervals\cite{efron1993} over 2{,}000 resamples; proportions in the human-criterion analysis are reported with Wilson intervals\cite{wilson1927} instead, because bootstrap intervals are degenerate or uninformative for observed proportions at the boundary. For macro-F1 the confusion matrix is resampled from a multinomial with the observed cell proportions, which is equivalent to resampling exchangeable items. The equivalence is verified against direct item resampling on every run; the check and its tolerance are documented in the released code.

Three null distributions are used for the trajectory-metric correlation, each of 10{,}000 draws, and each summarised by its mean, standard deviation, 2.5th and 97.5th percentiles and an empirical $p$-value
\begin{equation}
p = \frac{1 + \#\{r_{\mathrm{null}} \ge r_{\mathrm{obs}}\}}{1 + B},
\label{eq:pval}
\end{equation}
with $B=10{,}000$ draws. The three null comparisons are treated as one analysis family and are the only empirical $p$-values reported in this paper. The trajectory-null analysis itself was specified after the observed association had been examined and is therefore reported as an exploratory, mechanism-focused null analysis rather than as a confirmatory test of a prior hypothesis; the interval-exclusion statements in Section~\ref{subsec:remaining} are a separate, exploratory set and are not $p$-values. This is a one-sided upper-tail empirical $p$-value, because the question is whether the observed positive association is larger than the null rather than merely different from it. All three nulls are specified algorithmically below and are regenerable from the released analysis configuration, which fixes the random-number generator and its seeds.

\emph{Within-document permutation.} Each draw randomly reorders the observed sentence-level sentiment values within every document and recomputes both metrics. Each document's multiset of values is preserved exactly and only the ordering changes, so this null preserves more of the observed document-level structure than the synthetic nulls do. An association that does not exceed it cannot be attributed uniquely to the temporal ordering of those observed values.

\emph{Synthetic i.i.d.} For a document of $N$ sentences, $N$ independent values are drawn from $\mathcal{N}(0, \sigma^2)$ with $\sigma = 0.2488$, the pooled standard deviation of the sentence-level scores of the 1{,}288 trajectory-eligible relationship posts. Document lengths are matched to the observed distribution. The Gaussian is unbounded while the sentiment score is not, so this null matches the variance of the observed signal rather than its support; it is a secondary diagnostic and no inferential conclusion rests on it.

\emph{Synthetic bounded Gaussian random walk.} $s_1 = 0$ and $s_t = \mathrm{clip}(s_{t-1} + \varepsilon_t, -1, +1)$ with $\varepsilon_t \sim \mathcal{N}(0, \sigma_{\mathrm{step}}^2)$ independent across $t$ and $\sigma_{\mathrm{step}} = 0.2678$, the standard deviation of the absolute successive differences of those same 1{,}288 posts. This is a step-size proxy rather than an exact match to the variance of the signed increments, and is stated as such because the two quantities differ; no alternative step-size calibration was tested; lengths matched as above. The state bounds make this a bounded rather than a free random walk; the process is bounded and exhibits strong temporal dependence. It is reported to show the correlation such a process induces, not as a model of narrative.

\subsection{Preregistration and deviations}\label{subsec:prereg}

This study was \emph{not} preregistered. No external, timestamped registration exists for the hypotheses reported here, and none is claimed. That is stated plainly because the alternative -- describing analyses as confirmatory without a record -- cannot be verified by a reader, and is the practice Kerr\cite{kerr1998} named.

What does exist is an internal decision file recording 26 analysis choices, each written before the corresponding evaluation number was computed and each carrying a machine timestamp: the classification threshold, the label mappings for every corpus, the two-class treatment of ISEAR and EmpatheticDialogues, the trajectory scope of Section~\ref{subsec:corpora}, the exclusion of labMT from the baseline comparison, the choice of external criteria, and the reserved annotation-development subset. It is released with the analysis scripts. It is a within-project audit trail, not a preregistration, and it is offered as the former.

Table~\ref{tab:dev} classifies every analysis reported here. Analyses marked \emph{post hoc} were specified after some result had been observed and are reported as exploratory regardless of outcome.

\begin{table}[ht]
\centering
\small
\caption{Classification of the analyses reported in this paper. No analysis is described as confirmatory, because no external registration exists.}
\label{tab:dev}
\begin{tabular}{p{5.0cm}p{3.0cm}p{5.4cm}}
\toprule
Analysis & Classification & Basis \\
\midrule
Benchmark comparison against reference lexicons & Specified in advance & Threshold, mappings and baseline set fixed before scoring \\
Trajectory-metric permutation null & Post hoc & Specified after the association was observed; reported for the mechanism it tests, not as a test of a prior prediction \\
Watershed against a human criterion & Specified in advance & Annotation protocol fixed before any comparison \\
Annotator reliability & Specified in advance & Repeat items and integrity checks built into the protocol \\
Component variability of the composite metric & Post hoc & Substituted for a leave-one-out reading after the renormalisation artifact was identified \\
Abstention-aware evaluation & Post hoc & Specified after the zero-return behaviour was quantified \\
Length-stratified moderator analysis & Post hoc & Exploratory; only the longest stratum excludes zero \\
Archetype resolution by corpus & Post hoc & Descriptive; no criterion, no validation claim \\
\botrule
\end{tabular}
\end{table}

\subsection{Human annotation}\label{subsec:annotation}

The watershed metric is intended to locate the point at which a narrative's emotional direction changes. Testing that requires a human criterion. A correlation with another metric derived from the same signal cannot serve as one, for the reason established in Section~\ref{subsec:null}.

Annotators did not participate in the framework's development, were not told the study's hypotheses or which metric-level predictions were at stake, and were not shown any framework output, including predicted watershed positions. Each annotator worked from their own file and did not see any other annotator's responses. Annotators were shown posts from the reserved 200-post annotation-development subset with sentences numbered, and recorded the narrator's overall sentiment, the sentence number at which the emotional direction changes most clearly or zero if there is none, the overall arc, and a confidence rating on a five-point scale labelled in the interface as 1 = guess and 5 = certain. The interface accepted exactly one integer sentence index or zero for the turning point; ranges, multiple sentences and free-text qualifications were not permitted. Roughly one item in ten was shown twice, unmarked, to measure agreement with the annotator's own earlier judgement. Repeat items were drawn at random, without replacement, from those positions early enough in the sequence to allow a gap of at least 40 intervening items, and the second showing was placed at a random position beyond that gap. The tool recorded time on task and whether the post was scrolled to its end.

Four rounds were conducted. Two were rejected by automated integrity checks and are reported in Section~\ref{subsec:reliability}, because the rejections are themselves informative about the difficulty of the task. The third round, three independent annotators, passed all checks and is the source of the reliability estimates. A fourth, single-annotator round of 40 posts provides the preliminary comparison against the metric; it has no inter-rater reliability by construction and every figure derived from it is labelled preliminary.

\subsection{Ethics declarations}\label{subsec:ethics}

\subsubsection*{Human annotation}

The author is an independent researcher with no institutional affiliation and
therefore had no access to an institutional review board or research ethics
committee. No ethics approval was sought or obtained, and none is claimed. The
following is a full account of what the annotation study involved, so that the
editor and reviewers can judge it directly.

Four annotators took part across four rounds. Each was an adult volunteer known personally to the author, recruited informally and unpaid. No annotator stood in any supervisory, employment, academic-assessment or other dependency relationship to the author, and none had anything to gain or lose from the study's outcome. None participated in
the framework's development, none was told the study's hypotheses, and none was
shown any output of the framework. Each was told before beginning that they were
labelling text for a research study, that their responses would be analysed and
published, and that they could stop at any point; each agreed on that basis.
Participation was voluntary and could be withdrawn at any time.

The task was to read publicly posted text and record four judgements: an overall
sentiment, a sentence at which the emotional direction changes or an indication
that none does, an overall narrative arc, and a confidence rating on a five-point scale from 1, guess, to 5, certain. The annotation interface additionally recorded, for each item, an active-time measure, the number of times a response was revised, and whether the document had been scrolled to its end. Active time advances one second at a time and only while the browser tab is visible and the annotator has produced a keystroke, mouse movement, click or scroll within the preceding ninety seconds; total elapsed time from first display to submission is recorded separately, and every timing figure quoted in this paper is the active measure. These process measures were collected for quality control
and are reported in aggregate. No demographic, biometric, health, or otherwise
personal information was collected about any annotator, and no annotator is
identifiable in the released data, which records only an anonymous rater index.

In the author's assessment the study presents no more than minimal risk: it
involves adult volunteers performing a text-labelling task on publicly available
material, with no deception, no sensitive personal disclosure by the participants themselves, and no collection of identifying information. Annotators were nonetheless reading third-party accounts of relationship breakdown, infidelity and emotional distress, and could therefore encounter emotionally difficult material; they were told what the texts concerned before beginning and could stop at any point.
\subsubsection*{Third-party text}

The relationship corpus consists of posts made publicly on Reddit. Their authors
did not consent to research use and were not contacted. Several of the
subreddits sampled concern relationship breakdown, infidelity and emotional
distress, so although the posts are public they contain personal disclosures
made for support rather than for research. The following measures were taken.
Only post bodies were retained; usernames, comments and any other account
information were discarded at collection and are not held. The post text is not
redistributed: the release contains the collection script and a list of post
identifiers, so the corpus is partially reproducible via those identifiers for as long as the source material remains publicly accessible, and no copy of the text is published by us. Posts may be deleted or accounts removed, so exact reconstruction cannot be guaranteed. No quotation from any post appears in this paper, and no
individual post is described in a way that could identify its author. No attempt
was made to identify, contact, or infer attributes of any author.

The author notes that public availability is not the same as consent to
research use, and that this is an unresolved question for observational studies
on social media text generally rather than one settled by this paper.

\section{Results}\label{sec:results}

\subsection{What the instrument sees}\label{subsec:coverage}

The framework's 348-entry lexicon matches between 1.37\% and 2.95\% of tokens across corpora, against 7.00--14.85\% for VADER (Fig.~\ref{fig:scope}a, Table~\ref{tab:coverage}). On the relationship corpus it sees 28.4\% of what VADER sees.

Coverage is an important determinant of the empirical operating range. On documents of 1--9 words the framework returns a non-zero score for 7.2--23.8\% of items; on documents of 80--159 words, for 62.5--86.3\% (Fig.~\ref{fig:scope}b). At the corpus-average coverage rate of 2.48\%, a 395-word document -- the median length of this corpus -- would provide roughly ten matched tokens on average, making a non-zero score more plausible but not guaranteed, because matched terms can cancel; on a 12-word comment it corresponds on average to fewer than one matched word. Low coverage is survivable on long text but is often insufficient to produce a non-zero lexical signal on short text, and this is an important determinant of the accuracy figures in Section~\ref{subsec:benchmark}, alongside lexical polarity quality, the label mappings, the threshold and each corpus's class distribution.

\begin{figure}[ht]
\centering
\includegraphics[width=\linewidth]{fig1_coverage_scope.png}
\caption{Operating range of the instrument. (a) Fraction of tokens matched by the framework's 348-entry lexicon against VADER's 7{,}506-entry lexicon, per corpus. (b) Percentage of documents in each word-count band for which the framework returned a non-zero score; the relationship corpus is shown as a solid line. Bands containing fewer than 20 documents are omitted.}
\label{fig:scope}
\end{figure}

\begin{table}[ht]
\centering
\small
\caption{Lexicon coverage: fraction of tokens each system's lexicon matches. All lexicons were matched against the same tokenised representation of each corpus, so the figures differ only in lexicon content and not in preprocessing.}
\label{tab:coverage}
\begin{tabular}{lrrrrrrr}
\toprule
Corpus & PLEF & VADER & NRC & Empath & AFINN & Hu-Liu & PLEF as \% of VADER \\
\midrule
semeval2017  & 1.37\% & 7.00\%  & 9.39\%  & 23.82\% & 5.96\%  & 5.27\%  & 19.6\% \\
tweeteval    & 1.50\% & 7.65\%  & 10.33\% & 24.69\% & 6.71\%  & 6.01\%  & 19.6\% \\
meld         & 1.51\% & 14.85\% & 7.01\%  & 14.36\% & 11.28\% & 7.42\%  & 10.1\% \\
dailydialog  & 1.55\% & 11.23\% & 8.68\%  & 38.16\% & 9.56\%  & 7.47\%  & 13.8\% \\
relationship & 2.48\% & 8.74\%  & 8.66\%  & 34.81\% & 7.61\%  & 6.64\%  & 28.4\% \\
isear        & 2.54\% & 9.82\%  & 11.94\% & 41.66\% & 8.89\%  & 7.05\%  & 25.9\% \\
empathetic   & 2.76\% & 8.99\%  & 11.15\% & 49.64\% & 8.17\%  & 7.57\%  & 30.7\% \\
goemotions   & 2.95\% & 14.65\% & 11.52\% & 30.08\% & 13.26\% & 11.34\% & 20.1\% \\
\botrule
\end{tabular}
\end{table}

\subsection{Benchmark performance}\label{subsec:benchmark}

Under a common $\pm 0.05$ decision threshold, on the headline full-corpus macro-F1 comparison, the framework does not lead against any of the five reference systems on any evaluated corpus (Table~\ref{tab:f1}, Fig.~\ref{fig:bench}). The unweighted mean of the corpus-level macro-F1 values across the five three-class corpora is 0.406 for the framework against 0.508 for AFINN, 0.484 for Hu-Liu, 0.483 for VADER, 0.424 for Empath and 0.423 for NRC. Excluding TweetEval, which shares 42.8\% of its items with SemEval, changes the framework's mean to 0.403 and AFINN's to 0.497; the ordering is unaffected.

Two entries require comment. On MELD the framework scores 0.321 while a uniformly random three-class prediction scores 0.326, 95\% interval $[0.318, 0.333]$: on that corpus, where its coverage is lowest at 1.51\%, the framework's score falls within the simulated uniform-random reference distribution. No inferential test of the framework against that distribution is reported. The random benchmark is a Monte Carlo reference distribution: the mean macro-F1 over 1{,}000 independent assignments in which each item is given one of the three labels with equal probability, independently of its gold label. The interval is the 2.5th and 97.5th percentiles of those simulated values and is simulation-based rather than analytic. On DailyDialog it scores 0.420 against VADER's 0.350, which is its only apparent lead on a three-class corpus; Section~\ref{subsec:abstention} shows this to be largely driven by abstention on a corpus that is 84.9\% neutral, and the threshold-free macro-AUC on the same corpus is 0.584 against VADER's 0.731 (Fig.~\ref{fig:bench}b).

Under this positive/negative mean binary AUC over positive-vs-rest and negative-vs-rest, a threshold-free evaluation of the underlying scalar score whose values are not directly comparable with conventional three-class macro-AUC implementations, the framework is last on four of the five evaluated three-class corpora; on MELD it is fifth of six, ahead of Empath by 0.004 (0.541 against 0.537).

\begin{figure}[ht]
\centering
\includegraphics[width=\linewidth]{fig2_benchmark.png}
\caption{Benchmark performance against reference implementations. (a) macro-F1 with bootstrap 95\% confidence intervals over 2{,}000 resamples. The short dotted segment above each corpus marks that corpus's own random-prediction benchmark, computed separately for each because the expected macro-F1 of a uniform random classifier depends on the gold-label distribution: goemotions 0.327, semeval2017 0.319, meld 0.326, tweeteval 0.323, dailydialog 0.236. (b) positive/negative mean binary AUC over positive-vs-rest and negative-vs-rest, hereafter mean binary AUC over positive-vs-rest and negative-vs-rest, a threshold-free evaluation of the underlying scalar score rather than of any thresholded decision. Points, not bars, on an untruncated axis. This is not the conventional three-class macro-AUC and its values are not directly comparable with implementations that include neutral-versus-rest; the dotted line marks chance.}
\label{fig:bench}
\end{figure}

\begin{table}[ht]
\centering
\small
\caption{macro-F1 with bootstrap 95\% confidence intervals over 2{,}000 resamples. Leading zeros are omitted from interval bounds. The two-class corpora are reported separately and are not pooled with the three-class corpora.}
\label{tab:f1}
\begin{threeparttable}
\begin{tabular}{lccccc c}
\toprule
& \multicolumn{5}{c}{Three-class corpora} & \\
\cmidrule(lr){2-6}
System & goemotions & semeval2017 & meld\tnote{a} & tweeteval & dailydialog\tnote{b} & Mean \\
\midrule
PLEF   & .447 [.442,.451] & .423 [.418,.427] & .321 [.314,.329] & .421 [.417,.426] & .420 [.415,.424] & 0.406 \\
VADER  & .581 [.577,.585] & .541 [.537,.545] & .404 [.396,.413] & .541 [.537,.546] & .350 [.347,.353] & 0.483 \\
NRC    & .479 [.475,.483] & .450 [.445,.454] & .373 [.364,.381] & .456 [.452,.460] & .359 [.355,.362] & 0.423 \\
Empath & .469 [.465,.473] & .460 [.456,.465] & .335 [.327,.343] & .463 [.459,.467] & .392 [.389,.396] & 0.424 \\
AFINN  & \textbf{.604} [.600,.608] & \textbf{.556} [.552,.560] & \textbf{.431} [.423,.440] & \textbf{.556} [.552,.560] & .395 [.391,.398] & \textbf{0.508} \\
Hu-Liu & .555 [.550,.559] & .525 [.520,.529] & .405 [.396,.414] & .525 [.521,.530] & .411 [.407,.414] & 0.484 \\
\midrule
& \multicolumn{2}{c}{Two-class corpora} & & & & \\
\cmidrule(lr){2-3}
System & isear & empathetic & & & & \\
\midrule
PLEF   & .630 [.614,.645] & .564 [.556,.571] & & & & \\
VADER  & .645 [.632,.657] & \textbf{.724} [.718,.731] & & & & \\
NRC    & .585 [.573,.597] & .664 [.658,.671] & & & & \\
Empath & .585 [.573,.598] & .626 [.619,.633] & & & & \\
AFINN  & \textbf{.659} [.646,.672] & .686 [.680,.693] & & & & \\
Hu-Liu & .654 [.640,.668] & .669 [.663,.676] & & & & \\
\botrule
\end{tabular}
\begin{tablenotes}\footnotesize
\item[a] A uniformly random three-class prediction scores 0.326 on MELD (95\% interval $[0.318,0.333]$); the framework scores 0.321.
\item[b] The framework's lead here is largely driven by abstention; see Section~\ref{subsec:abstention}.
\item[c] Unweighted mean of the five corpus-level macro-F1 values; corpora are not weighted by size.
\end{tablenotes}
\end{threeparttable}
\end{table}

\subsection{Abstention is not a neutral prediction}\label{subsec:abstention}

When the lexicon matches nothing the framework returns exactly zero, which the threshold maps to ``neutral''. Mapping zero to neutral conflates absence of lexical evidence with a neutral prediction, and can inflate measured performance on neutral-heavy corpora. The headline scores in Table~\ref{tab:f1} apply that mapping, as every reported benchmark does; they are the ordinary thresholded scores and not an abstention-aware decision rule. The analysis below separates the two readings but does not replace the benchmark with a corrected metric. Separating the two changes the interpretation of Table~\ref{tab:f1} in two places (Table~\ref{tab:abst}, Fig.~\ref{fig:gase}b).

The framework produces a non-zero score on 12.1--27.8\% of items in the three-class corpora. On those items it remains below VADER by 0.064--0.080 macro-F1, so the restricted operating range does not by itself account for the accuracy gap: the framework is also less accurate than VADER on the fired subsets of these three-class corpora. That subset is defined by the framework's own behaviour, so it is not a neutral sample of the corpus and this is a post-selection comparison.

DailyDialog is the informative case. The framework's headline 0.420 there falls to 0.249 on the 15.2\% of items where it fires. Its apparent lead is largely driven by abstention on 84.8\% of a corpus in which, after the processing described in Table~\ref{tab:reconcile}, 84.9\% of the 101{,}026 retained utterances carry the no-emotion label; the abstentions are therefore scored as correct neutral predictions. This proportion is a property of our processed utterance-level subset and is close to, but not identical with, the figure usually quoted for the canonical release.

On the two-class corpora the direction reverses. Where the framework fires it exceeds VADER by 0.039 on ISEAR (0.752 against 0.712) and by 0.017 on EmpatheticDialogues (0.828 against 0.812). The framework exceeded VADER on the fired subset in these two corpora; these are the only evaluated fired subsets on which it does so. The differences are point estimates and no test of their significance is reported; both are higher-firing subsets of the two-class corpora.

\begin{table}[ht]
\centering
\caption{Abstention-aware evaluation. ``Fires'' is the percentage of documents for which the lexicon matched at least one term. $\Delta$ is computed from unrounded scores and rounded independently of the two columns it compares, so it may differ by one unit in the last place from their displayed difference.}
\label{tab:abst}
\begin{tabular}{lrrrrrr}
\toprule
Corpus & Fires & PLEF (all) & VADER (all) & PLEF (fired) & VADER (fired) & $\Delta$ \\
\midrule
goemotions   & 27.8\% & 0.447 & 0.581 & 0.474 & 0.549 & $-0.075$ \\
semeval2017  & 23.2\% & 0.423 & 0.541 & 0.457 & 0.527 & $-0.070$ \\
meld         & 12.1\% & 0.321 & 0.404 & 0.371 & 0.447 & $-0.076$ \\
tweeteval    & 23.0\% & 0.421 & 0.541 & 0.444 & 0.524 & $-0.080$ \\
dailydialog  & 15.2\% & 0.420 & 0.350 & 0.249 & 0.313 & $-0.064$ \\
\midrule
isear        & 39.2\% & 0.630 & 0.645 & \textbf{0.752} & 0.712 & $+0.039$ \\
empathetic   & 36.7\% & 0.564 & 0.724 & \textbf{0.828} & 0.812 & $+0.017$ \\
\botrule
\end{tabular}
\end{table}

\subsection{The trajectory-metric correlation does not exceed its null}\label{subsec:null}

On the 1{,}288 relationship-related posts with at least four sentences, the correlation between the arc asymmetry and the watershed drop is $r_{\mathrm{obs}} = +0.598$, 95\% CI $[+0.547, +0.646]$. Taken alone, an association of this size between two independently named metrics would ordinarily be read as convergent validity for both.

It is not unusually large relative to the within-document permutation distribution (Table~\ref{tab:nulls}, Fig.~\ref{fig:null}a). Over 10{,}000 permutation replicates of the full 1{,}288-document analysis the null mean is $+0.6074$ with standard deviation $0.0264$ and 95\% interval $[+0.5544, +0.6583]$; $6{,}425$ of the $10{,}000$ draws are at least as large as the observed value, giving an empirical $p = 0.643$ by Equation~\eqref{eq:pval}, with a Monte Carlo standard error of 0.005. Against the synthetic i.i.d.\ null the observed value gives $p = 0.440$. Against the bounded Gaussian random-walk null, whose mean is $+0.8497$, all 10{,}000 draws exceed the observed value; the tabulated $p=1.000$ therefore records that no draw fell below the observed correlation, not a population probability of one.

The observed association is therefore not unusually large relative to what the same sentiment values produce when their order is destroyed, and is far below what a bounded Gaussian random walk of comparable step size produces.

This does not depend on the arc window. The sensitivity comparison was run at 1{,}000 draws per rule rather than the 10{,}000 of the primary analysis, so its piecewise $p$ differs in the third decimal from the $p=0.643$ reported above for the same specification. Repeating the whole analysis with $k=\lfloor N/4 \rfloor$ for every document gives an observed $r=+0.611$ against a null mean of $+0.614$ ($p=0.560$); with $k=\lfloor N/3 \rfloor$ for every document, $r=+0.635$ against $+0.633$ ($p=0.473$); and with the piecewise rule used here, $r=+0.598$ against $+0.607$ ($p=0.630$), against $p=0.643$ at 10{,}000 draws. In all three specifications the observed value is not unusually large relative to its corresponding null distribution. Changing the window shifts the observed and null correlations within a similar range while the inferential conclusion is unchanged, which is the behaviour expected if the association is generated by the signal the two metrics share rather than by the ordering of the sentences.

Equations~\eqref{eq:lewi} and~\eqref{eq:nava} suggest a mechanism: both quantities are differences between the mean of an early portion of $S$ and the mean of a late portion, and for short sequences with the split near the middle they are close to the same quantity. The algebra alone does not establish how large the induced correlation should be; the permutation result does, and it is the empirical test rather than the algebraic resemblance that carries the conclusion. In this dataset the observed association does not exceed the association generated by the specified within-document order-randomisation null, and it therefore does not provide evidence of convergent validity for either metric in this evaluation.

This case illustrates a broader methodological caution. The empirical evidence here comes from a single corpus of 1{,}288 relationship posts, and the claim is about that evidence. Pairs of trajectory statistics derived from a common per-sentence signal can correlate strongly through the signal they share, and a within-document permutation, which tests that one mechanism rather than all structural mechanisms, is inexpensive to compute and provides a direct test of whether an observed association exceeds the association induced by randomising sentence order while holding each document's sentence-level values fixed. We suggest it be considered routinely.

\begin{table}[ht]
\centering
\caption{The arc-watershed correlation against three null distributions, on 1{,}288 relationship-related posts. Each null is 10{,}000 draws. $p$ is the empirical proportion of null draws at least as large as the observed correlation, by Equation~\eqref{eq:pval}.}
\label{tab:nulls}
\begin{threeparttable}
\begin{tabular}{lrrlr}
\toprule
& mean $r$ & SD & 95\% interval & $p$ \\
\midrule
Observed ($n=1{,}288$) & $+0.598$ & --- & $[+0.547, +0.646]$\tnote{a} & --- \\
\midrule
Within-document permutation & $+0.6074$ & $0.0264$ & $[+0.5544, +0.6583]$ & $0.643$ \\
Synthetic i.i.d.\ noise & $+0.5945$ & $0.0231$ & $[+0.5486, +0.6388]$ & $0.440$ \\
Synthetic bounded random walk & $+0.8497$ & $0.0091$ & $[+0.8314, +0.8671]$ & $1.000$ \\
\botrule
\end{tabular}
\begin{tablenotes}\footnotesize
\item[a] Bootstrap interval on the observed correlation, 2{,}000 resamples; the null intervals are percentiles of the 10{,}000 null draws.
\end{tablenotes}
\end{threeparttable}
\end{table}

\begin{figure}[ht]
\centering
\includegraphics[width=\linewidth]{fig3_nulls_moderator.png}
\caption{(a) The observed arc-watershed correlation (vertical line) against three null distributions of 10{,}000 draws each. The permutation null preserves each document's sentence-level values and destroys only their order; the two synthetic nulls are specified in Section~\ref{subsec:nullmethod}. Histograms are the 10{,}000 actual simulated draws per null, not a parametric approximation; the draws are released as \texttt{null\_draws.csv}. (b) The pronoun-power association within the relationship corpus, by document length, with bootstrap 95\% confidence intervals. Only the longest stratum excludes zero.}
\label{fig:null}
\end{figure}

\subsection{The watershed metric against a human criterion}\label{subsec:human}

An annotator independent of the framework's development and blinded to its outputs and to the study's hypotheses worked through 40 presented items from the reserved annotation-development subset, spending a median of 124 seconds of active time per item. Table~\ref{tab:annacct} gives the full accounting from presented items to the comparisons that follow, because the intermediate steps are what make the terminal counts interpretable.

Of the 35 posts on which both a human answer and a defined watershed index exist, the annotator judged that 16 contain no turning point. The metric returned an index on all 35, because the watershed procedure as implemented always selects one index from a non-empty candidate set and has no explicit no-turning-point state. The current watershed implementation has no explicit no-turning-point state and cannot represent the absence of a turn, the annotator assigned no turning point to 46\% of these posts, a rate observed under one annotator and one protocol rather than an estimate of prevalence. This limitation follows from the defined watershed procedure rather than from the observed sample, and this structural limitation is directly relevant to interpreting the comparison, because the annotation analysis indicates that, in this sample, disagreement was concentrated in the existence judgement.

On the remaining 19 posts, where the annotator did identify a turning point, the metric selected the same sentence on none of them (Table~\ref{tab:human}). The chance benchmark must respect the support of Equation~\eqref{eq:lewi}: the metric can return only the $N_i-2$ interior positions $2,\dots,N_i-1$, not all $N_i$ sentences. The exact-match benchmark is therefore the mean over the 19 posts of $1/(N_i-2)$, and the within-one-sentence benchmark is the mean over posts of $|\{h_i-1, h_i, h_i+1\} \cap \{2,\dots,N_i-1\}| / (N_i-2)$, computed for each post from its own sentence count $N_i$ and the annotator's selection $h_i$ rather than approximated as $3/N_i$. A further correction is required: the annotation instructions placed no restriction on which sentence a human could select, so a human selection of sentence 1 or sentence $N$ lies outside the metric's support and its exact-match probability is zero rather than $1/(N_i-2)$. The benchmark therefore contributes $1/(N_i-2)$ for interior selections and zero for boundary selections. On the 19 comparison posts, whose sentence counts range from 7 to 114 with a median of 30, one of the 19 selections falls on a boundary (sentence 17 of a 17-sentence post) and therefore contributes zero, giving 4.67\% for exact match; the within-one-sentence benchmark, which already intersects the annotator's neighbourhood with the admissible support, is 14.2\%. The naive $1/N_i$ and $3/N_i$ forms give 4.3\% and 12.9\%. The per-post values behind both benchmarks -- each post's sentence count, the annotator's selection, whether it lies on the boundary, and its exact contribution -- are released as \texttt{chance\_per\_post.csv}. The observed agreement of 0.0\% exact and 10.5\% within one sentence is below both corrected benchmarks, but neither difference is evidence of below-chance performance. Under the same per-post chance model the expected number of exact matches across the 19 posts is 0.89, so observing none has probability 0.40 and is the single most likely outcome; the expected number of within-one matches is 2.70 against two observed, with probability 0.48 of two or fewer. The comparison therefore fails to demonstrate agreement with the annotator rather than demonstrating disagreement, and at $n=19$ it has almost no power to separate the two. The observed exact-match and within-one-sentence agreement rates are both below their corresponding support-aware chance benchmarks.

For the arc metric, agreement with the annotator was 37.1\% against a majority-class baseline of 42.9\%, with Cohen's $\kappa = +0.020$. The distributions of the two are close -- the total variation distance is 0.043 -- while item-level agreement is essentially chance-level. A metric that reproduces the marginal distribution can appear successful in aggregate summaries while having little item-level agreement. Reporting the distribution of predicted arcs, without item-level agreement against a criterion, cannot distinguish the two cases.

These figures come from a single annotator and are preliminary. They are reported with the reliability ceiling of Section~\ref{subsec:reliability} beside them.

\begin{table}[ht]
\centering
\small
\caption{Accounting for the human annotation, from presented items to the comparisons reported in Table~\ref{tab:human}. Every count is reconstructed from the annotation records.}
\label{tab:annacct}
\begin{tabular}{lrl}
\toprule
Stage & $n$ & Note \\
\midrule
Items presented to the annotator & 40 & includes hidden repeats \\
\quad of which repeat showings & 4 & used only for intra-annotator agreement \\
Unique posts & 36 & first showing of each \\
Posts matched to a scored document & 36 & \\
\quad of which the watershed index is defined & 35 & one post has fewer than four sentences \\
\quad\quad annotator marked \emph{no} turning point & 16 & metric returned an index on all 35 \\
\quad\quad annotator marked a turning point & 19 & the comparison set for exact agreement \\
\botrule
\end{tabular}
\end{table}

\begin{table}[ht]
\centering
\small
\caption{The trajectory metrics against a human criterion. Single annotator; preliminary.}
\label{tab:human}
\begin{threeparttable}
\begin{tabular}{lrr}
\toprule
Measure & Value & Chance or baseline \\
\midrule
Watershed, exact sentence match ($n=19$) & 0.0\% & 4.67\%\tnote{a} \\
Watershed, within one sentence ($n=19$) & 10.5\% & 14.2\%\tnote{a} \\
Posts where the annotator marked no turning point & 16 of 35 & metric returned an index on all 35 \\
Arc agreement ($n=35$) & 37.1\% & 42.9\% (majority class) \\
Arc, Cohen's $\kappa$ & $+0.020$ & 0 \\
\botrule
\end{tabular}
\begin{tablenotes}\footnotesize
\item[a] Computed on the metric's admissible support $\{2,\dots,N_i-1\}$ from each post's own sentence count. One of the 19 annotator selections is a boundary sentence the metric cannot return and contributes zero. The naive $1/N_i$ and $3/N_i$ forms give 4.3\% and 12.9\%. The per-post values behind both benchmarks -- each post's sentence count, the annotator's selection, whether it lies on the boundary, and its exact contribution -- are released as \texttt{chance\_per\_post.csv}.
\end{tablenotes}
\end{threeparttable}
\end{table}

\subsection{The turning-point judgement is not reliably reproducible in this study}\label{subsec:reliability}

Whether a metric can be validated against a criterion depends on whether the criterion is stable. Three annotation rounds were conducted and the answer is that, for the turning-point judgement, it is not.

In the third round, three annotators each labelled 100 posts under a designed overlap: 174 distinct posts in total, of which 26 were labelled by all three, 74 by exactly one, and the remainder by exactly two. The pairwise shared sets are 53, 51 and 48 posts, each containing the same 26 three-way items plus 27, 25 and 22 pair-only items respectively. Krippendorff's $\alpha$ was 0.177 for sentiment and 0.002 for arc on the 26 posts labelled by all three, against 0.67, which we adopted a priori as the minimum level for tentative criterion-level use in this study, following the discussion of reliability thresholds in Krippendorff\cite{krippendorff2004}.

The turning-point judgement requires two agreement measures, which differ and are easily conflated. There was \emph{no unanimous} three-way agreement on the turning-point sentence on any of the 26 posts labelled by all three. \emph{Pairwise} agreement must be reported as two quantities, because the response space contains a null category alongside the sentence positions and a single percentage mixes them. Counting a shared ``no turning point'' as agreement, raw pairwise agreement was 15.1\% (8 of 53), 29.4\% (15 of 51) and 25.0\% (12 of 48); but that figure cannot be compared against a chance rate defined over interior sentences alone, because the two are not the same response space. The two questions are therefore separated. \emph{Existence} agreement asks whether both annotators agreed that a turn exists at all, and is compared against the chance rate implied by each annotator's own rate of recording no turning point. \emph{Location} agreement asks, conditional on both having selected a sentence, whether they selected the same one, and is compared against the mean of $1/(N_i-2)$ over exactly those posts. The two behave very differently. \emph{Existence} agreement across the three pairs was 58.5\%, 62.7\% and 52.1\% against marginal-based chance rates of 65.0\%, 54.4\% and 52.1\%: at or below chance on two of the three pairs. \emph{Location} agreement, conditional on both annotators having selected a sentence, was 23.3\% (7 of 30), 34.6\% (9 of 26) and 27.8\% (5 of 18) against chance rates of 7.8\%, 5.6\% and 8.2\%: approximately three- to six-fold above the corresponding support-based chance rates, in a small sample.

The direction of this result is the opposite of what a single mixed percentage suggests. In this sample, disagreement was concentrated in whether a turn existed rather than in its location conditional on both annotators having selected one. The conditional location comparisons rest on 30, 26 and 18 posts respectively, so this decomposition is suggestive rather than established. The corresponding chance rate is computed per pair as the mean over that pair's own shared posts of $1/(N_i-2)$, the admissible support of Equation~\eqref{eq:lewi}, rather than from a single median sentence count. Unanimity never occurred on any of the 26 posts labelled by all three. Both figures are reported because the unanimous one alone would overstate the case.

Using the unmarked repeated items, agreement with the annotator's own earlier judgement was 71--75\% for sentiment and 27--88\% for arc, but 18--29\% for the turning point. In the fourth, single-annotator round, agreement with the same annotator's own earlier answer on repeated posts was 75\% for sentiment, 100\% for arc and 0\% for the turning point, on four repeats; these four are descriptive only and estimate nothing.

In this annotation study the reported reliability coefficients were below the adopted threshold for sentiment and arc. The turning-point judgement decomposes: annotators agreed at or below chance on whether a turn exists, and approximately three- to six-fold above chance on where it is when both said one existed, in a small sample. Which judgement is weakest therefore depends on the measure, and the four quantities are not on a common scale. On raw repeat agreement the turning point is lowest, at 18--29\% against 71--75\% for sentiment. On chance-corrected agreement the arc is lowest, with $\alpha = 0.002$, indistinguishable from chance, while pairwise location agreement runs approximately three- to six-fold above its chance rate in a small sample, and existence agreement sits at or below chance. No single judgement is weakest on every measure, and we do not claim one is. No gold standard was constructed from the three-annotator round, and the reason is not that a majority vote would have failed to resolve items. A two-of-three majority resolves 18 of the 26 posts annotated by all three, 69.2\%. The primary obstacle is empirical: of those 18, eight are resolved as ``no turning point'' and only ten name a specific sentence, which is too small a set for a meaningful sentence-selection validation analysis. Secondarily, no unanimous agreement occurred, within-annotator repeat agreement was 18--29\%, and the sentiment and arc coefficients on the same items are 0.177 and 0.002, so these characteristics make the resulting majority criterion too uncertain for sentence-level validation in this sample. The counts are reported in the supplementary annotation audit.

Two earlier rounds were rejected by automated checks before this one. The checks and their rejection thresholds were written into the annotation tooling before the responses of the round they rejected were inspected, and are listed in the supplementary material; the rejections are therefore decisions against a stated rule rather than post-hoc judgements about unwelcome data. In the first, the recorded turning-point response was identical between every pair of annotators on all 200 shared items. The comparison is between the raw responses, so two annotators who both recorded ``no turning point'' count as identical alongside those who selected the same sentence. Under an independent-response null over the metric's admissible sentence positions, computed from that round's own posts as the mean of $1/(N_i-2)$, the expected pairwise rate of selecting the same sentence is 8.44\%; the observed 100\% is far above that even before the shared-zero responses are considered, and the check that rejected the round is defined on the raw responses rather than on that benchmark. In the second, across 600 pairwise comparisons -- three annotator pairs over 200 items -- sentiment, arc and confidence either all three matched or all three differed. Not one comparison showed exactly one or exactly two matching fields. The rejection rests on that observed pattern rather than on a computed chance rate: three fields judged independently would be expected to produce partial matches, and a complete absence of them across 600 comparisons is not a pattern independent judgement produces. Both patterns are highly inconsistent with that stated independent-response null; the exact expected rates are given in the supplementary material. The rejection rules, their thresholds, the per-round statistics and the raw response files for both rejected rounds are given in the supplementary material; the checks are part of the released tooling. Both rejections are reported because a study that discards two annotation rounds should say so.

\subsection{Composite metric components}\label{subsec:gase}

The composite relationship-health metric is a weighted sum of a sentiment term, a normalised emotion entropy, an attachment term and a conflict-marker term, with weights 0.30, 0.25, 0.25 and 0.20. It is assessed by leave-one-out ablation with the remaining weights renormalised.

The informative quantity is the variability of each component (Fig.~\ref{fig:gase}a, Table~\ref{tab:gase}). The attachment term has a standard deviation below 0.05 on seven of eight corpora and the conflict term on all eight. Together they carry 45\% of the composite weight while varying little within those corpora, so on them the composite is close to a fixed offset plus the two remaining terms. Only on the relationship corpus does the attachment standard deviation exceed 0.05, at 0.127, which is consistent with the corpus for which it was designed.

The leave-one-out deltas themselves are reported in the supplementary material with a caveat that applies to any renormalised ablation: because the removed component's weight is redistributed, the delta measures a component's value relative to the others rather than its informativeness. Under the renormalised leave-one-out calculation, removing the sentiment term increases the composite, because its mean is near 0.01 while the other three terms average close to 1.0. The standard deviations describe within-corpus variation; they do not by themselves establish whether a component is informative, since a component with little marginal variation could still contribute conditionally.

\begin{table}[ht]
\centering
\caption{Standard deviation of each composite component. Values below 0.05 are treated here as indicating little within-corpus variation.}
\label{tab:gase}
\begin{tabular}{lrrrr}
\toprule
Corpus & sd(sentiment) & sd(entropy) & sd(attachment) & sd(conflict) \\
\midrule
dailydialog  & 0.102 & 0.234 & 0.011 & 0.026 \\
empathetic   & 0.203 & 0.371 & 0.010 & 0.018 \\
goemotions   & 0.180 & 0.355 & 0.010 & 0.044 \\
isear        & 0.228 & 0.402 & 0.017 & 0.036 \\
meld         & 0.118 & 0.239 & 0.007 & 0.028 \\
relationship & 0.091 & 0.227 & \textbf{0.127} & 0.023 \\
semeval2017  & 0.139 & 0.327 & 0.011 & 0.027 \\
tweeteval    & 0.149 & 0.327 & 0.011 & 0.031 \\
\botrule
\end{tabular}
\end{table}

\begin{figure}[ht]
\centering
\includegraphics[width=\linewidth]{fig4_gase_abstention.png}
\caption{(a) Standard deviation of each composite component per corpus; the dotted line marks 0.05. The attachment and conflict terms vary below that level on seven and eight corpora respectively. (b) Difference in macro-F1 between the framework and VADER on the documents where the framework's lexicon fired, against the proportion of documents on which it fired. The framework exceeds VADER only on the two two-class corpora.}
\label{fig:gase}
\end{figure}

\subsection{Remaining metric relationships}\label{subsec:remaining}

The remaining metric relationships were examined on the relationship corpus with bootstrap intervals. These are exploratory analyses and no multiplicity adjustment was applied. The number of associations, strata and alternative definitions examined here raises the false-discovery risk accordingly, and no confirmatory inference is drawn from any of them. Association is reported as association and is not treated as evidence of construct validity; the interval-exclusion column below is descriptive and should not be read as a family of hypothesis tests and, where possible, against externally sourced criteria not derived from the framework's own output (Table~\ref{tab:corr}).

The pronoun-power association is $r=+0.379$ when the power index retains the shared pronoun term and $r=+0.090$ when it does not, a reduction of 76\%; the large reduction after removing the shared pronoun term indicates that a substantial portion of the association under the first definition is attributable to the two indices sharing their inputs. Against the dominance dimension of the NRC valence-arousal-dominance lexicon\cite{mohammad2018vad}, which is externally sourced and is not constructed from the pronoun counts that enter either index, the association is $r=+0.021$ with an interval spanning zero. The document-level dominance value is the unweighted mean of the lexicon's dominance rating over those tokens of the document that appear in the lexicon, using the same tokenisation as every other measure here; tokens absent from the lexicon are omitted rather than imputed as zero, and no frequency weighting or rescaling is applied. A document matching no lexicon entry would have an undefined value and be excluded pairwise from the correlation; no such document occurs, and all 1{,}308 posts contribute.

The volatility metric was tested against a sentence-level intensity sequence derived from the NRC Emotion Intensity lexicon\cite{mohammad2018intensity}, which is not derived from the framework's sentiment core. That sequence is constructed as follows: for each sentence, the intensity value is the unweighted mean over the sentence's tokens that appear in the lexicon of the maximum intensity assigned to that token across the lexicon's emotion categories, with absent tokens omitted and a sentence containing no matched token assigned zero. The same mean absolute successive difference of Section~\ref{subsec:framework} is then applied to that sequence, with no additional normalisation, so the two volatility measures differ only in the lexicon that produces the per-sentence values. The association is $r=-0.096$ with an interval excluding zero. The observed association is negative, contrary to the predicted positive direction; it is reported as measured and no theoretical interpretation is offered for its sign. This is an intermediate criterion rather than the definitive one: validation against an instrument-based affect-lability measure administered to the writers themselves cannot be performed using the present public-corpus design and remains future work.

The sentiment--distortion association is $r=+0.018$ with an interval spanning zero.

One relationship is length-dependent. Within the relationship corpus the pronoun-power association increases with document length, from $r=+0.014$ in the 20--149 word stratum to $r=+0.161$ in the 600-word-and-above stratum, and only the longest stratum has an interval excluding zero (Fig.~\ref{fig:null}b). The same relationship is absent across corpora, $r=+0.032$ over seven corpus-level points, so the effect is within-corpus and not a general property of length.

\begin{table}[ht]
\centering
\small
\caption{Metric relationships on the relationship corpus ($n=1{,}308$), with bootstrap 95\% confidence intervals.}
\label{tab:corr}
\begin{tabular}{llr}
\toprule
Relationship & $r$ [95\% CI] & Interval excludes zero \\
\midrule
Pronoun index $\times$ power index, earlier definition & $+0.379$ \ [$+0.332, +0.422$] & yes \\
Pronoun index $\times$ power index, shared term removed & $+0.090$ \ [$+0.045, +0.137$] & yes \\
Power index $\times$ NRC-VAD dominance (external) & $+0.021$ \ [$-0.025, +0.069$] & no \\
Volatility $\times$ NRC intensity volatility (intermediate lexical criterion) & $-0.096$ \ [$-0.166, -0.030$] & yes, wrong sign \\
Sentiment $\times$ distortion index & $+0.018$ \ [$-0.048, +0.077$] & no \\
Diagnostic: composite $\times$ semantic-fusion metric (internal overlap) & $+0.108$ \ [$+0.048, +0.164$] & yes \\
\botrule
\end{tabular}
\end{table}

The last row is a diagnostic rather than a hypothesis test. The composite metric carries 35\% of the weight in the semantic-fusion metric by definition, yet correlates with it at $r=+0.108$, which follows from the near-constancy of two of its components reported in Section~\ref{subsec:gase}.

\subsection{Implementation and reproducibility}\label{subsec:repro}

The framework core is standard-library only, verified by a purity test that imports it in a fresh process with third-party packages blocked. Output is deterministic: a test scores an identical input set in two separate processes and compares the serialised results, which are byte-for-byte identical. The analysis pipeline was configured not to construct a consensus gold standard below the prespecified reliability threshold, so that an unreliable criterion could not enter a downstream comparison unnoticed. That is a study-design rule, not a property of the software. The evaluation harness additionally requires the reference baseline packages, NumPy, SciPy and Matplotlib; the two dependency sets are separate and this is stated rather than described as a dependency-free system.

All analyses were run under CPython 3.11.9 on 64-bit Windows with eight logical cores. The evaluation harness requires vaderSentiment, afinn, empath, nltk, labMTsimple, NumPy, SciPy and Matplotlib; exact versions are pinned in the released environment file, and the NRC lexicons are obtained separately under their research licence. Throughput on eight cores, summed over the eight evaluation corpora, is 5{,}501{,}892 words in 61.3 minutes of scoring time, a rate of 1{,}496 words per second. This covers the twelve metric families and the derived semantic-fusion diagnostic, which is computed in the same pass. Per-corpus rates range from 765 to 1{,}939 words per second and vary with document length, so a single-corpus rate should not be combined with the total elapsed time.

The watershed index is sensitive to its own definition. A variant that searches for sign changes in the first derivative, rather than the maximum absolute second difference of Equation~\eqref{eq:lewi}, selects the same sentence on 71.4\% of documents in the relationship corpus; where the two differ, the median separation between the selected sentences is 5 and the median absolute difference in the reported drop is 0.073. Implementations of a watershed metric should therefore state the exact rule used, since plausible alternatives disagree on 28.6\% of documents in this corpus.

\section{Discussion}\label{sec:discussion}

\subsection{What is and is not supported}

The original study formulation contained seven hypotheses concerning the framework's metrics. Since no external registration exists (Section~\ref{subsec:prereg}), those labels are retired here and each underlying claim is re-evaluated on its own terms; Table~\ref{tab:scoreboard} reports the outcomes, and the system-comparison claim is subsumed by the benchmark results of Section~\ref{subsec:benchmark}. None of the original validity hypotheses receives positive support in this analysis. The outcomes are not equivalent in kind: some are contradicted by an interpretable test, others are simply not demonstrated because the criterion is unreliable or the sample too small, and the distinction is stated with each entry rather than collapsed into a single verdict. The trajectory-metric association does not exceed its null; the watershed metric does not match a human criterion and cannot represent the absence of a turning point; the volatility metric correlates with an externally sourced intensity measure in the direction opposite to the prediction; the original pronoun-power association is substantially attributable to shared inputs and is close to zero against an external criterion, with an interval spanning zero; two of four composite components are near-constant; and the sentiment--distortion association is null.

Two empirical findings are positive under their stated conditions; an exploratory archetype diagnostic is also reported in the supplementary material but is not validated. On the two two-class corpora, restricted to documents where its lexicon fires, the framework exceeds VADER by 0.039 and 0.017 macro-F1. The pronoun-power association increases with document length within the relationship corpus. An exploratory archetype metric is reported in the supplementary material. It has no external criterion, no validation is claimed for it, and no conclusion in this paper rests on it.

\begin{table}[ht]
\centering
\small
\caption{Historical study hypotheses and current evidence. These are the seven hypotheses of the original study formulation. These are historical study hypotheses, not preregistered confirmatory hypotheses; see Section~\ref{subsec:prereg}.}
\label{tab:scoreboard}
\begin{tabular}{p{4.1cm}p{2.6cm}p{6.6cm}}
\toprule
Hypothesis & Outcome & Evidence \\
\midrule
Watershed matches human turning points & Not supported (preliminary) & no exact matches in $n=19$ against an expected 0.89; $P(0)=0.40$ under the chance model, so this fails to demonstrate agreement rather than demonstrating disagreement; single annotator; the procedure cannot represent absence \\
Pronoun index predicts power asymmetry & Not supported & $r=+0.090$; external criterion $r=+0.021$, interval spans zero \\
Volatility tracks emotional volatility & Not supported & $r=-0.096$ against an intermediate lexical criterion; negative, contrary to the predicted direction \\
Composite components contribute independently & Not supported & Two of four components have sd $<0.05$ on 7--8 of 8 corpora \\
Arc classification matches human arcs & Not supported (preliminary) & 37.1\% against a 42.9\% majority baseline, $\kappa=+0.020$ \\
Sentiment predicts cognitive distortion & Not supported & $r=+0.018$, interval spans zero \\
Arc $\times$ watershed as convergent validity & Withdrawn & $+0.598$ observed against a $+0.607$ permutation null mean \\
\botrule
\end{tabular}
\end{table}

\subsection{Implications beyond this framework}

Three methodological implications extend beyond this particular implementation.

Pairs of trajectory statistics derived from a common per-sentence signal can correlate strongly through the signal they share. The correlation reported here is closely reproduced when the sentence order that the metrics are supposed to measure is destroyed. A within-document order randomisation is cheap to compute and provides a direct test of whether an observed association exceeds the association induced by randomising sentence order while holding each document's sentence-level values fixed; without it, a structural identity is difficult to distinguish from convergent validity. We suggest that such a null be considered when reporting associations between trajectory statistics.

Lexicon-based instruments can exhibit a coverage floor. A 348-entry lexicon produces a non-zero signal on most documents near the median length of the relationship corpus and fails to produce one on most short-form documents. The relationship between coverage, document length and the ability to score at all is measurable in advance, costs nothing, and empirically characterises the range in which a given implementation frequently produces a lexical signal. Accuracy on corpora outside that range is strongly influenced by the thresholded handling of non-firing cases and should not be read solely as evidence about lexical sentiment discrimination.

Returning zero because no lexical evidence was found is not a prediction of neutrality, and conflating the two inflates apparent accuracy on neutral-heavy corpora. The DailyDialog result is an instance: an apparent lead of 0.070 macro-F1 over VADER becomes a deficit of 0.064 once abstentions are separated from predictions, and under the threshold-free mean binary AUC over positive-vs-rest and negative-vs-rest used here it leads no system on any evaluated three-class corpus.

\subsection{Construct validity of narrative turning points}

The reliability results have a consequence beyond this framework. In the three-annotator reliability round, turning-point judgements never coincided unanimously on any jointly annotated post; in the separate fourth round, a single annotator spent a median of 124 seconds of active time per post. The reported inter-annotator reliability coefficients for sentiment and arc were below the threshold adopted here. The turning-point judgement decomposes into two questions that behave oppositely: annotators agreed at or below chance on whether a turn exists, and approximately three- to six-fold above chance on where it is once both said one existed, in a small sample.

The pattern is more specific than uniformly random judgement, and it runs opposite to what a single mixed agreement figure implies. Pairs of annotators located the turn approximately three- to six-fold above the support-based chance rate whenever both agreed one existed, in a small sample, while agreeing at or below chance on whether one existed; the three never agreed unanimously, and within-annotator repeat agreement was 18--29\%. These results are consistent with conditional localisation being more stable than the existence judgement in this sample. One possible explanation is that readers differ in how much emotional change they require before calling it a turn, rather than in where they see the change occur; that could produce this pattern of chance-level existence agreement with well-above-chance localisation. This study does not distinguish that explanation from other sources of annotation disagreement, including instruction wording, sentence segmentation, and annotator training or fatigue. If that holds more widely, it constrains what any watershed metric can be validated against. In a direct observed-score validation comparison of this kind, a criterion with very poor reliability cannot provide strong evidence for or against validity. The reported $\alpha$ values of 0.177 for sentiment and 0.002 for arc are both below the threshold adopted here; turning-point reliability was assessed separately, because the response is a sentence selection with a possible null category rather than a categorical label, and under those measures it showed no unanimous three-way agreement and 18--29\% within-annotator repeat agreement. No single coefficient describes the reliability of all three judgements, and none of them supports a criterion-based validation claim in this study. Whether the judgement becomes reliable with training, a narrower definition, or a different genre is an open question this study cannot answer; it is measured here rather than assumed in either direction.

\subsection{Limitations}\label{subsec:limitations}

The human comparison rests on 36 unique posts presented as 40 item appearances, including four hidden repeats, annotated by one person; it is preliminary throughout. The reliability estimates rest on 26 three-way overlapping posts. Both are small, and the intervals reported reflect that.

A pooled relationship-domain classification score is not estimable from the labels available. We did not identify labels of sufficient reliability for the relationship corpus, the only other relationship-adjacent labelled corpus is two-class and cannot be pooled with three-class corpora, and constructing an aggregate from these would require averaging incommensurable quantities. An exploratory, point-estimate development-set comparison is therefore the only one available, on the 36 posts carrying human sentiment labels (16 positive, 15 negative, 5 neutral). On those posts the framework scores macro-F1 $0.338$, against $0.365$ for VADER, $0.340$ for NRC, $0.407$ for Empath, $0.363$ for AFINN and $0.402$ for Hu-Liu: it has the lowest macro-F1 among the six systems compared. These are point estimates only. Given $n=36$ we do not interpret them as evidence of a relationship-domain performance ranking; they are reported with their denominator rather than as a domain estimate. No relationship-domain aggregate is reported, because none can be computed from labels of sufficient reliability.

The framework's sentiment lexicon was not expanded during this study, because modifying it after observing the results would constitute tuning on the evaluation data. The coverage limitation reported here is therefore a property of the instrument as specified, and whether a larger lexicon would remove it is untested.

The relationship corpus retains no post-level timestamps, so no temporal analysis of the corpus is possible and the collection window cannot be reconstructed from the released data.

The composite and archetype metrics are exploratory. The archetype metric has no external criterion; the supplementary material reports where it returns output, not whether that output is correct.

Corpus-level pooling is affected by the 42.8\% overlap between SemEval-2017 and TweetEval; means are reported both ways and the ordering is unchanged.

\subsection{What the framework is for}

An advantage of interpretable instruments is that their outputs can be inspected and contested rather than only measured against a benchmark\cite{doshivelez2017}; that same property is what made the failures reported here diagnosable. The instrument that survives this evaluation is narrow and its empirical operating boundaries are characterised in these corpora. The framework core is deterministic and uses only the Python standard library; it produces per-sentence attributable output and was evaluated on relationship-related posts with a median length of 394.5 words. It is not competitive with reference lexicon implementations at sentiment classification and is not offered as such. Its trajectory metrics did not demonstrate evidence of narrative-shape validity in this evaluation.

The value of the exercise is in the negative results and the methods that produced them, and in the demonstration that a set of validity claims of this kind can be substantially clarified by three inexpensive checks: a permutation null, a reliability estimate for the human criterion, and a comparison against reference rather than reimplemented baselines.

\subsection{Conclusion}

This study tested a twelve-metric lexicon-based framework on 309{,}833 corpus
records and one collected corpus of 1{,}308 relationship-related Reddit posts.
What it establishes is narrow. The association between the framework's arc and
watershed metrics on 1{,}288 posts is not unusually large relative to a
within-document order-randomisation null, so it is not evidence of convergent
validity for either metric in this evaluation. The lexicon matches 1.37--2.95\%
of tokens, which restricts the document lengths on which the instrument produces
any signal at all, and its one apparent classification lead follows from mapping
absence of lexical evidence to a neutral prediction.

What it does not establish is equally important. It does not show that
trajectory metrics are invalid in general: the trajectory evidence comes from a
single purposively sampled corpus, and the order-randomisation null tests one
mechanism rather than every structural explanation. It does not show that the
watershed metric disagrees with human judgement: the direct comparison rests on
19 posts and one annotator, where the expected number of chance matches is below
one, and it therefore fails to demonstrate agreement without demonstrating
disagreement. It does not establish that the turning-point construct is
unreliable in general, only that it was not reliably reproducible under this
protocol.

The framework is currently fit for describing lexical coverage and abstention
behaviour on long-form text, and for the trajectory computations on documents of
at least four sentences where those quantities are wanted descriptively. It is
not fit for criterion-level use.

The single most important next experiment is a larger annotation study that
separates the existence of a narrative turn from its location. The decomposition
reported here, on 30, 26 and 18 conditional pairwise items, suggests that
annotators disagree mainly about whether a turn occurs rather than where it is.
If that holds at scale it would change how such metrics should be defined and
validated, and it is testable with a protocol that admits a null response, an
interval, and multiple candidate turns.

\section*{Author contributions}

R.K.M. designed the study, implemented the framework and the analysis, conducted
the evaluation and wrote the manuscript.

\section*{Data availability}

The analysis scripts, the rebuilt implementation, the decision file recording all
analysis choices with timestamps, the annotation tooling including its integrity
checks, and the manifests with input checksums are available at
\url{https://github.com/rahul-k-mishra/plef}. The public corpora are obtainable
from their original sources, cited in Section~\ref{subsec:corpora}. The NRC
lexicons are available from their author under a research licence and are not
redistributed. The relationship corpus consists of public Reddit posts; the
collection scripts are provided and the post text is not redistributed. The
annotation responses for all rounds, including the two rejected ones, are
included. The implementation and results of the previous version of this work, which this
evaluation corrects, are preserved at the \texttt{pre-revision} tag of the same
repository.

\section*{Additional information}

\textbf{Competing interests} The author declares no competing interests.

\section*{Funding}

This research received no external funding.

\section*{Ethics statement}

This study used publicly available Reddit posts. Only post bodies were
retained; usernames, comments and all other account information were discarded
at collection, no personally identifiable information was retained, the post
text is not redistributed, and no post is quoted or described in a way that
could identify its author.

The study also involved four adult volunteers who annotated that text. The
author is an independent researcher with no institutional affiliation and had no
access to an institutional review board or research ethics committee; no ethics
approval was sought or obtained, and none is claimed. Each annotator was told
before beginning that they were labelling text for a research study, that their
responses would be analysed and published, and that they could stop at any
point, and each agreed on that basis. No demographic, biometric, health or other
personal information was collected about any annotator, and none is identifiable
in the released data, which records only an anonymous rater index. No clinical
claim is made and no clinical data were used. A full account is given in
Section~\ref{subsec:ethics}.

\begin{thebibliography}{99}

\bibitem{hutto2014}
Hutto, C.\ J.\ \& Gilbert, E.
VADER: a parsimonious rule-based model for sentiment analysis of social media text.
\emph{Proc.\ 8th Int.\ AAAI Conf.\ Weblogs and Social Media} 216--225 (2014).

\bibitem{mohammad2013}
Mohammad, S.\ M.\ \& Turney, P.\ D.
Crowdsourcing a word-emotion association lexicon.
\emph{Comput.\ Intell.} \textbf{29}, 436--465 (2013).

\bibitem{fast2016}
Fast, E., Chen, B.\ \& Bernstein, M.\ S.
Empath: understanding topic signals in large-scale text.
\emph{Proc.\ CHI Conf.\ Human Factors in Computing Systems} 4647--4657 (2016).

\bibitem{nielsen2011}
Nielsen, F.\ {\AA}.
A new ANEW: evaluation of a word list for sentiment analysis in microblogs.
\emph{Proc.\ ESWC Workshop on Making Sense of Microposts} 93--98 (2011).

\bibitem{hu2004}
Hu, M.\ \& Liu, B.
Mining and summarizing customer reviews.
\emph{Proc.\ 10th ACM SIGKDD Int.\ Conf.\ Knowledge Discovery and Data Mining} 168--177 (2004).

\bibitem{dodds2010}
Dodds, P.\ S.\ \& Danforth, C.\ M.
Measuring the happiness of large-scale written expression: songs, blogs, and presidents.
\emph{J.\ Happiness Stud.} \textbf{11}(4), 441--456 (2010).

\bibitem{reagan2016}
Reagan, A.\ J., Mitchell, L., Kiley, D., Danforth, C.\ M.\ \& Dodds, P.\ S.
The emotional arcs of stories are dominated by six basic shapes.
\emph{EPJ Data Sci.} \textbf{5}, 31 (2016).

\bibitem{mohammad2018vad}
Mohammad, S.\ M.
Obtaining reliable human ratings of valence, arousal, and dominance for 20,000 English words.
\emph{Proc.\ 56th Annu.\ Meeting Assoc.\ Comput.\ Linguistics} 174--184 (2018).

\bibitem{mohammad2018intensity}
Mohammad, S.\ M.
Word affect intensities.
\emph{Proc.\ 11th Int.\ Conf.\ Language Resources and Evaluation (LREC 2018)}
(Miyazaki, Japan, 2018).

\bibitem{demszky2020}
Demszky, D.\ et al.
GoEmotions: a dataset of fine-grained emotions.
\emph{Proc.\ 58th Annu.\ Meeting Assoc.\ Comput.\ Linguistics} 4040--4054 (2020).

\bibitem{rashkin2019}
Rashkin, H., Smith, E.\ M., Li, M.\ \& Boureau, Y.-L.
Towards empathetic open-domain conversation models: a new benchmark and dataset.
\emph{Proc.\ 57th Annu.\ Meeting Assoc.\ Comput.\ Linguistics} 5370--5381 (2019).

\bibitem{li2017}
Li, Y.\ et al.
DailyDialog: a manually labelled multi-turn dialogue dataset.
\emph{Proc.\ 8th Int.\ Joint Conf.\ Natural Language Processing} 986--995 (2017).

\bibitem{poria2019}
Poria, S.\ et al.
MELD: a multimodal multi-party dataset for emotion recognition in conversations.
\emph{Proc.\ 57th Annu.\ Meeting Assoc.\ Comput.\ Linguistics} 527--536 (2019).

\bibitem{rosenthal2017}
Rosenthal, S., Farra, N.\ \& Nakov, P.
SemEval-2017 Task 4: sentiment analysis in Twitter.
\emph{Proc.\ 11th Int.\ Workshop on Semantic Evaluation} 502--518 (2017).

\bibitem{barbieri2020}
Barbieri, F., Camacho-Collados, J., Espinosa Anke, L.\ \& Neves, L.
TweetEval: unified benchmark and comparative evaluation for tweet classification.
\emph{Findings of EMNLP} 1644--1650 (2020).

\bibitem{scherer1994}
Scherer, K.\ R.\ \& Wallbott, H.\ G.
Evidence for universality and cultural variation of differential emotion response patterning.
\emph{J.\ Pers.\ Soc.\ Psychol.} \textbf{66}, 310--328 (1994).

\bibitem{kuppens2010}
Kuppens, P., Allen, N.\ B.\ \& Sheeber, L.\ B.
Emotional inertia and psychological maladjustment.
\emph{Psychol.\ Sci.} \textbf{21}, 984--991 (2010).

\bibitem{houben2015}
Houben, M., Van Den Noortgate, W.\ \& Kuppens, P.
The relation between short-term emotion dynamics and psychological well-being: a meta-analysis.
\emph{Psychol.\ Bull.} \textbf{141}, 901--930 (2015).

\bibitem{deerwester1990}
Deerwester, S., Dumais, S.\ T., Furnas, G.\ W., Landauer, T.\ K.\ \& Harshman, R.
Indexing by latent semantic analysis.
\emph{J.\ Am.\ Soc.\ Inf.\ Sci.} \textbf{41}, 391--407 (1990).

\bibitem{landauer1998}
Landauer, T.\ K., Foltz, P.\ W.\ \& Laham, D.
An introduction to latent semantic analysis.
\emph{Discourse Process.} \textbf{25}, 259--284 (1998).

\bibitem{tausczik2010}
Tausczik, Y.\ R.\ \& Pennebaker, J.\ W.
The psychological meaning of words: LIWC and computerized text analysis methods.
\emph{J.\ Lang.\ Soc.\ Psychol.} \textbf{29}, 24--54 (2010).

\bibitem{boyd2022}
Boyd, R.\ L., Ashokkumar, A., Seraj, S.\ \& Pennebaker, J.\ W.
\emph{The Development and Psychometric Properties of LIWC-22}
(University of Texas at Austin, 2022).

\bibitem{krippendorff2004}
Krippendorff, K.
\emph{Content Analysis: An Introduction to Its Methodology}, 2nd edn
(Sage, 2004).

\bibitem{cohen1960}
Cohen, J.
A coefficient of agreement for nominal scales.
\emph{Educ.\ Psychol.\ Meas.} \textbf{20}, 37--46 (1960).

\bibitem{efron1993}
Efron, B.\ \& Tibshirani, R.\ J.
\emph{An Introduction to the Bootstrap}
(Chapman \& Hall, 1993).

\bibitem{wilson1927}
Wilson, E.\ B.
Probable inference, the law of succession, and statistical inference.
\emph{J.\ Am.\ Stat.\ Assoc.} \textbf{22}, 209--212 (1927).

\bibitem{kerr1998}
Kerr, N.\ L.
HARKing: hypothesizing after the results are known.
\emph{Pers.\ Soc.\ Psychol.\ Rev.} \textbf{2}, 196--217 (1998).

\bibitem{mcadams2013}
McAdams, D.\ P.\ \& McLean, K.\ C.
Narrative identity.
\emph{Curr.\ Dir.\ Psychol.\ Sci.} \textbf{22}, 233--238 (2013).

\bibitem{doshivelez2017}
Doshi-Velez, F.\ \& Kim, B.
Towards a rigorous science of interpretable machine learning.
Preprint at \url{https://arxiv.org/abs/1702.08608} (2017).

\end{thebibliography}



\end{document}
